\section{Operating Framework: The (Y, M, X) Triple}\label{sec:2}

This section defines the unit of analytical work that the iteration roster
operates on. We adopt the same definitions used internally throughout the
project's specs, plans, and memory artifacts so that the supervisor's
comments map directly onto the documents that govern subsequent iterations.

\subsection{Definitions}

The triple \((Y, M, X)\) consists of:

\begin{itemize}
  \item \(Y\) — the welfare-relevant outcome variable on which a target
    population's exposure to the wage-to-capital transition is measured.
    \(Y\) is operationalised per population. Examples used in closed and
    active iterations include the Section~J information-and-communication
    employment share for young Colombian workers, the realised volatility
    of a household consumption basket, and the cross-sectional differential
    between productive-capital returns and a wage-indexed consumption
    basket.
  \item \(X\) — the macro risk hypothesised to block the wage-to-capital
    transition for the target population. \(X\) identification is empirical
    and population-specific. The first-cut iteration question is always:
    which macro variable, at which lag, currently kills wage earners'
    attempts to enter productive-capital ownership in this population? \(X\)
    enters the regression as a lag-distributed right-hand-side variable;
    the lag structure is pinned in the spec text before the panel is
    constructed.
  \item \(M\) — the settlement geometry of the candidate hedging
    instrument. \(M\) is introduced for completeness; it is declared out of
    scope for empirical validation and deferred to Stage~2 of each
    iteration. The on-chain operational details that determine the
    feasibility of \(M\) for a given \((Y, X)\) pair are surveyed in
    Section~\ref{sec:4} for the supervisor's reference; the rest of the
    memorandum does not require familiarity with them.
\end{itemize}

\subsection{Stage discipline}

Each iteration proceeds through three stages with explicit exit criteria:

\begin{enumerate}
  \item \textbf{Empirical risk validation.} Estimate
    \(\hat{\beta}_{\mathrm{composite}}\) of \(Y\) on lagged \(X\) under a
    pre-registered specification. Exit on a positive-\(\hat{\beta}\)
    verdict at the pre-committed significance threshold. We treat this
    stage as the object of the supervisor's review.
  \item \textbf{Ideal-scenario M sketch.} Conditional on a Stage-1 PASS,
    propose an instrument geometry on the on-chain substrate that
    \emph{would} settle the empirical \(\hat{\beta}_{\mathrm{composite}}\) if
    deployed. Stage~2 does not require live counterparty capital. Exit on
    a settlement-feasible construction that hedges the validated
    correlation.
  \item \textbf{Live deployment.} Conditional on Stage~2 and live
    counterparty capital, deploy and run an execution test. This stage is
    out of scope for the present memorandum.
\end{enumerate}

The stage discipline is enforced by an internal anti-fishing rule: stage
drift, in particular allowing Stage~2 instrument-design considerations to
re-enter the Stage-1 specification choice, is treated as a silent design
adjustment under the project's pre-registration discipline (see
Section~\ref{sec:5} and Section~\ref{sec:7.4}). The empirical exercises
the supervisor is being asked to review are entirely Stage-1.

\subsection{Iteration order}

The default iteration order is target-population dominant. We fix the
target population first; we fix \(Y\) on a candidate
inequality- or wage-to-capital-transition exposure measure for that
population; we enumerate \(X\) candidates from the empirical risk surface
that the population faces; for each surviving \(X\) we search the
on-chain-eligible \(M\) space for tradability. The pair \((Y, X)\) only
graduates to instrument design once a settlement-feasible \(M\) exists. A
closed iteration with verdict FAIL or \texttt{EXIT} informs the
\(X\)-search prior for the next population. We do \emph{not} re-run the
same \((Y, X)\) at adjusted thresholds after a FAIL; the within-iteration
anti-fishing invariants carry forward to subsequent iteration design.

\subsection{Acknowledged limitations}

Two limitations of the framework are load-bearing for the supervisor's
review and worth stating explicitly in this section rather than in
Section~\ref{sec:7}.

First, the iteration order described above does not, by itself, address
cross-iteration multiple-comparison concerns. Iteration selection (which
populations and which \(X\) candidates enter the roster) is a project-level
decision that is not pre-registered in the same way that within-iteration
specifications are pre-registered. We treat this as an open methodological
question (Section~\ref{sec:7.5}).

Second, identification at Stage~1 is correlation-identified rather than
causally identified. The project's PASS-verdict iteration (Pair~D)
explicitly hedges the correlation rather than asserting a causal channel
\citep[per the closure memo
\texttt{memory/project\_pair\_d\_phase2\_pass.md}, Reality-Checker
FLAG~\#1]{}. We discuss the identification posture and its implications for
moving from Stage~1 to Stage~2 in Section~\ref{sec:7.1}.
